% GENERATED FILE. Edit canonical lessons, then run npm run generate:books.
% canonical-source-hash: fnv1a64:71ac37ee171320ab
% canonical-lessons: RU-C427-gavan, RU-C427-mayak, RU-C427-matros, RU-C427-kapitan, RU-C427-tonut

\chapter{In the harbour}
\label{ch:ru-in-the-harbour}
\hlchaptermodality{427}

\begin{chapteropening}
\textbf{By the end of this chapter:} \emph{I can talk about a harbour: the harbour, the lighthouse, the sailor, the captain and a boat that sinks.}

Complete the last lesson of chapter 427: I can talk about a harbour: the harbour, the lighthouse, the sailor, the captain and a boat that sinks.
\end{chapteropening}

\section[\texorpdfstring{\ru{гавань} (gávan)}{гавань (gávan)}]{\ru{гавань} (gávan) — a harbour}

\label{lesson:RU-C427-gavan}



\begin{quote}
Before the new one: say the Russian for an oar, then the Russian for to row.
\end{quote}

\subsection*{You'll want to know: \ru{гавань}}
\textbf{\ru{гавань}} (\emph{gávan}) — \textquotedblleft{}a harbour\textquotedblright{}.

\begin{quote}
\ru{Вечером} \ru{в} \ru{гавань} \ru{вернулись} \ru{все} \ru{лодки}.
\end{quote}

\begin{quote}
\emph{Vécherom v gávan vernúlis vse lódki.}
\end{quote}

\begin{quote}
In the evening all the boats came back to the harbour.
\end{quote}

Feminine, though it ends in -\ru{ь}.

\begin{cousinweb}
\emph{\ru{гавань}} was borrowed from Dutch \emph{haven} 'a harbour' (EtymDB), from Middle Dutch \emph{havene}.
\end{cousinweb}

\begin{grammarlens}[title={What you've built}]
One more word for telling, explaining and describing at B1.
\end{grammarlens}

\subsection*{Your turn}
\begin{itemize}
\raggedright
  \item \emph{Say it:} \emph{gávan}
  \item \emph{Say it:} \emph{gávan}, once more
  \item \emph{Say it:} a sentence of your own with \emph{gávan}
  \item \emph{Recall:} say the Russian for an oar, then the Russian for to row, then say \emph{gávan} again
\end{itemize}

\begin{tcolorbox}[breakable,colback=teal!4,colframe=teal!35!black,title={Before you move on}]
What does \ru{гавань} mean? (\textquotedblleft{}A harbour\textquotedblright{}.) Say it once more.
\end{tcolorbox}
\section[\texorpdfstring{\ru{маяк} (mayák)}{маяк (mayák)}]{\ru{маяк} (mayák) — a lighthouse; a beacon}

\label{lesson:RU-C427-mayak}



\begin{quote}
Before the new one: say the Russian for to row, then the Russian for a harbour.
\end{quote}

\subsection*{You'll want to know: \ru{маяк}}
\textbf{\ru{маяк}} (\emph{mayák}) — \textquotedblleft{}a lighthouse; a beacon\textquotedblright{}.

\begin{quote}
\ru{На} \ru{скале} \ru{стоит} \ru{старый} \ru{маяк}.
\end{quote}

\begin{quote}
\emph{Na skalé stoít stáryy mayák.}
\end{quote}

\begin{quote}
An old lighthouse stands on the rock.
\end{quote}

Masculine; the stress moves: \emph{\ru{маяка}} (\emph{mayaká}).

\begin{cousinweb}
\emph{\ru{маяк}} is usually traced to a Turkic word for a pole or a marker. EtymDB lists the Russian word with the sense 'lighthouse' but no source, so the line is hedged.
\end{cousinweb}

\begin{grammarlens}[title={What you've built}]
One more word for telling, explaining and describing at B1.
\end{grammarlens}

\subsection*{Your turn}
\begin{itemize}
\raggedright
  \item \emph{Say it:} \emph{mayák}
  \item \emph{Say it:} \emph{mayák}, once more
  \item \emph{Say it:} a sentence of your own with \emph{mayák}
  \item \emph{Recall:} say the Russian for to row, then the Russian for a harbour, then say \emph{mayák} again
\end{itemize}

\begin{tcolorbox}[breakable,colback=teal!4,colframe=teal!35!black,title={Before you move on}]
What does \ru{маяк} mean? (\textquotedblleft{}A lighthouse; a beacon\textquotedblright{}.) Say it once more.
\end{tcolorbox}
\section[\texorpdfstring{\ru{матрос} (matrós)}{матрос (matrós)}]{\ru{матрос} (matrós) — a sailor, a seaman}

\label{lesson:RU-C427-matros}



\begin{quote}
Before the new one: say the Russian for a harbour, then the Russian for a lighthouse.
\end{quote}

\subsection*{You'll want to know: \ru{матрос}}
\textbf{\ru{матрос}} (\emph{matrós}) — \textquotedblleft{}a sailor, a seaman\textquotedblright{}.

\begin{quote}
\ru{Старый} \ru{матрос} \ru{рассказывал} \ru{нам} \ru{о} \ru{море}.
\end{quote}

\begin{quote}
\emph{Stáryy matrós rasskázyval nam o móre.}
\end{quote}

\begin{quote}
The old sailor told us about the sea.
\end{quote}

Masculine. \emph{\ru{Моряк}} is anyone who works at sea; \emph{\ru{матрос}} is a deckhand or a junior rank.

\begin{cousinweb}
\emph{\ru{матрос}} was borrowed from Dutch \emph{matroos} 'a sailor, a seaman', which EtymDB traces to French \emph{matelot}.
\end{cousinweb}

\begin{grammarlens}[title={What you've built}]
One more word for telling, explaining and describing at B1.
\end{grammarlens}

\subsection*{Your turn}
\begin{itemize}
\raggedright
  \item \emph{Say it:} \emph{matrós}
  \item \emph{Say it:} \emph{matrós}, once more
  \item \emph{Say it:} a sentence of your own with \emph{matrós}
  \item \emph{Recall:} say the Russian for a harbour, then the Russian for a lighthouse, then say \emph{matrós} again
\end{itemize}

\begin{tcolorbox}[breakable,colback=teal!4,colframe=teal!35!black,title={Before you move on}]
What does \ru{матрос} mean? (\textquotedblleft{}A sailor, a seaman\textquotedblright{}.) Say it once more.
\end{tcolorbox}
\section[\texorpdfstring{\ru{капитан} (kapitán)}{капитан (kapitán)}]{\ru{капитан} (kapitán) — a captain}

\label{lesson:RU-C427-kapitan}



\begin{quote}
Before the new one: say the Russian for a lighthouse, then the Russian for a sailor.
\end{quote}

\subsection*{You'll want to know: \ru{капитан}}
\textbf{\ru{капитан}} (\emph{kapitán}) — \textquotedblleft{}a captain\textquotedblright{}.

\begin{quote}
\ru{Капитан} \ru{сказал}, \ru{что} \ru{море} \ru{спокойное}.
\end{quote}

\begin{quote}
\emph{Kapitán skazál, chto móre spokóynoye.}
\end{quote}

\begin{quote}
The captain said the sea was calm.
\end{quote}

Masculine. Also the captain of a team.

\begin{cousinweb}
\emph{\ru{капитан}} is usually traced through Italian \emph{capitano} or French \emph{capitaine} to Late Latin \emph{capitāneus} 'a chief', from \emph{caput} 'head' (EtymDB lists those words; the Russian entry has no source).
\end{cousinweb}

\begin{grammarlens}[title={What you've built}]
One more word for telling, explaining and describing at B1.
\end{grammarlens}

\subsection*{Your turn}
\begin{itemize}
\raggedright
  \item \emph{Say it:} \emph{kapitán}
  \item \emph{Say it:} \emph{kapitán}, once more
  \item \emph{Say it:} a sentence of your own with \emph{kapitán}
  \item \emph{Recall:} say the Russian for a lighthouse, then the Russian for a sailor, then say \emph{kapitán} again
\end{itemize}

\begin{tcolorbox}[breakable,colback=teal!4,colframe=teal!35!black,title={Before you move on}]
What does \ru{капитан} mean? (\textquotedblleft{}A captain\textquotedblright{}.) Say it once more.
\end{tcolorbox}
\section[\texorpdfstring{\ru{тонуть} (tonút)}{тонуть (tonút)}]{\ru{тонуть} (tonút) — to sink; to drown}

\label{lesson:RU-C427-tonut}



\begin{quote}
Before the new one: say the Russian for a sailor, then the Russian for a captain.
\end{quote}

\subsection*{You'll want to know: \ru{тонуть}}
\textbf{\ru{тонуть}} (\emph{tonút}) — \textquotedblleft{}to sink; to drown\textquotedblright{}.

\begin{quote}
\ru{Старая} \ru{лодка} \ru{медленно} \ru{тонет}.
\end{quote}

\begin{quote}
\emph{Stáraya lódka médlenno tónet.}
\end{quote}

\begin{quote}
The old boat is slowly sinking.
\end{quote}

Imperfective; perfective \emph{\ru{утонуть}} (of a person) or \emph{\ru{затонуть}} (of a ship). Present \emph{\ru{тону}, \ru{тонешь}}.

\begin{cousinweb}
\emph{\ru{тонуть}} is usually traced to the same root as \emph{\ru{топить}} 'to drown, to sink', which EtymDB takes from Proto-Slavic \emph{topiti}. EtymDB lists \emph{\ru{тонуть}} itself without a source.
\end{cousinweb}

\begin{grammarlens}[title={What you've built}]
One more word for telling, explaining and describing at B1.
\end{grammarlens}

\subsection*{Your turn}
\begin{itemize}
\raggedright
  \item \emph{Say it:} \emph{tonút}
  \item \emph{Say it:} \emph{tonút}, once more
  \item \emph{Say it:} a sentence of your own with \emph{tonút}
  \item \emph{Recall:} say the Russian for a sailor, then the Russian for a captain, then say \emph{tonút} again
\end{itemize}

\begin{tcolorbox}[breakable,colback=teal!4,colframe=teal!35!black,title={Before you move on}]
What does \ru{тонуть} mean? (\textquotedblleft{}To sink; to drown\textquotedblright{}.) Say it once more.
\end{tcolorbox}
